\documentclass[12pt]{article}
\usepackage[margin=0.85in]{geometry}
\usepackage{enumitem}
\usepackage{listings}
\usepackage{xcolor}
\usepackage{framed}
\usepackage{amssymb}
\usepackage{needspace}

\lstset{
  basicstyle=\ttfamily\small,
  breaklines=true,
  frame=single,
  backgroundcolor=\color{gray!8}
}

\title{CS449 Lecture Handout: Input and Files}
\author{Dr.\ Jack Lange}
\date{Lecture 2}

\begin{document}
\maketitle

\section*{tar}

\Needspace{0.22\textheight}
\noindent\textbf{What tar is for:}
\vspace{0.14\textheight}

\Needspace{0.16\textheight}
\noindent\textbf{Example:} bundle this week's lab files \quad\texttt{tar -czf} / \texttt{tar -xzf}
\vspace{0.10\textheight}

\section*{C is parsed top down}

\Needspace{0.30\textheight}
\noindent\textbf{What ``top down'' changes when you write a .c file:}
\vspace{0.22\textheight}

\Needspace{0.30\textheight}
\noindent\textbf{Compare with Java:}
\vspace{0.22\textheight}

\Needspace{0.28\textheight}
\begin{lstlisting}
int main(void) {
    foo();
}
void foo(void) { }
\end{lstlisting}
\vspace{0.12\textheight}

\newpage
\section*{Types of input}

\noindent Command line \hspace{3.5em} Console (\texttt{scanf}) \hspace{3.5em} Streams

\Needspace{0.22\textheight}
\noindent\textbf{argv --- what it is:}
\vspace{0.14\textheight}

\Needspace{0.14\textheight}
\noindent\textbf{argv[0] is:}
\vspace{0.10\textheight}

\Needspace{0.14\textheight}
\noindent\textbf{Example:} run \texttt{./cmdargs hello 42}
\vspace{0.08\textheight}

\Needspace{0.40\textheight}
\noindent\textbf{scanf --- format string, \texttt{\&}, and the return value:}
\vspace{0.26\textheight}

\Needspace{0.24\textheight}
\noindent\textbf{Suppose} the call is
\begin{lstlisting}
int n;
int matched = scanf("%d", /* */);
\end{lstlisting}
\vspace{0.12\textheight}

\newpage
\section*{Streams}

\Needspace{0.24\textheight}
\noindent\textbf{What a stream is (FIFO):}
\vspace{0.16\textheight}

\Needspace{0.34\textheight}
\noindent\textbf{stdin / stdout / stderr:}
\vspace{0.24\textheight}

\Needspace{0.32\textheight}
\noindent\textbf{printf vs fprintf:}
\vspace{0.14\textheight}

\begin{lstlisting}
int printf(const char *format, ...);
int fprintf(FILE *stream, const char *format, ...);
\end{lstlisting}

\newpage
\Needspace{0.55\textheight}
\noindent\textbf{Picture:} program, OS buffer, console --- mark a flush
\begin{framed}
\begin{minipage}[t][0.44\textheight]{\textwidth}
\vspace{0.4em}
\end{minipage}
\end{framed}

\newpage
\section*{The shell and streams}

\Needspace{0.32\textheight}
\noindent\textbf{Why the shell can mix files and streams:}
\vspace{0.22\textheight}

\Needspace{0.36\textheight}
\noindent\textbf{What each does:}
\begin{itemize}[leftmargin=*]
  \item \texttt{>} \quad \texttt{ls > files.out}
  \item \texttt{>>} \quad \texttt{ls >> all\_files.out}
  \item \texttt{<} \quad \texttt{wc < files.out}
  \item \texttt{|} \quad \texttt{ls | grep .c | wc}
\end{itemize}
\vspace{0.22\textheight}

\newpage
\section*{Special streams}

\Needspace{0.70\textheight}
\noindent\textbf{/dev/null:}
\vspace{0.14\textheight}

\noindent\textbf{/dev/zero:}
\vspace{0.14\textheight}

\noindent\textbf{/dev/random:}
\vspace{0.14\textheight}

\noindent\textbf{Question:} How is \texttt{/dev/zero} different from \texttt{/dev/null}?
\vspace{0.22\textheight}

\newpage
\section*{File I/O in programs}

\Needspace{0.24\textheight}
\noindent\textbf{Who manages console streams vs files you open:}
\vspace{0.14\textheight}

\Needspace{0.36\textheight}
\noindent\textbf{fopen / fclose:}
\vspace{0.16\textheight}

\begin{lstlisting}
FILE *file = fopen(name, /* mode */);
fclose(file);
\end{lstlisting}
\vspace{0.12\textheight}

\newpage
\section*{ASCII vs raw bytes}

\Needspace{0.32\textheight}
\noindent\textbf{Character \texttt{"7"} vs integer 7:}
\vspace{0.22\textheight}

\Needspace{0.28\textheight}
\noindent\textbf{Example:}
\begin{lstlisting}
int i = 7;
fprintf("%d\n", i);
\end{lstlisting}
\noindent Memory: \hspace{8em} File:
\vspace{0.12\textheight}

\noindent\textbf{Example:}
\begin{lstlisting}
int i = -7;
fprintf("%d\n", i);
\end{lstlisting}
\noindent Memory: \hspace{8em} File:
\vspace{0.12\textheight}

\newpage
\Needspace{0.55\textheight}
\noindent\textbf{Picture:} memory on the left, file bytes on the right
\begin{framed}
\begin{minipage}[t][0.44\textheight]{\textwidth}
\vspace{0.4em}
\end{minipage}
\end{framed}

\newpage
\section*{Explicit file I/O}

\Needspace{0.24\textheight}
\noindent\textbf{Why data has to be in memory first:}
\vspace{0.14\textheight}

\Needspace{0.32\textheight}
\noindent\textbf{explicit vs implicit access:}
\vspace{0.22\textheight}

\Needspace{0.36\textheight}
\noindent\textbf{fread / fwrite:}
\vspace{0.22\textheight}

\begin{lstlisting}
int x = 5;
FILE *outfile = fopen(argv[1], "wb");
fwrite(&x, sizeof(x), 1, outfile);
\end{lstlisting}
\noindent memory: \hspace{8em} file:
\vspace{0.12\textheight}

\newpage
\section*{File position}

\Needspace{0.32\textheight}
\noindent\textbf{Who remembers the file offset:}
\vspace{0.22\textheight}

\Needspace{0.24\textheight}
\noindent\textbf{Suppose} you read 5 bytes, then read again.
\vspace{0.14\textheight}

\Needspace{0.20\textheight}
\noindent\textbf{Example:} \texttt{fseek(fstream, 16, SEEK\_SET)}
\vspace{0.12\textheight}

\newpage
\Needspace{0.55\textheight}
\noindent\textbf{Picture:} file as a row of bytes --- before a 5-byte read, after it, after that \texttt{fseek}
\begin{framed}
\begin{minipage}[t][0.44\textheight]{\textwidth}
\vspace{0.4em}
\end{minipage}
\end{framed}

\vspace{2em}
\begin{center}
\textit{This handout may not cover every topic from the source slides.}
\end{center}

\end{document}
